% Application description for the FitBuddy PIA (verified against source at commit 28d1adb, v3.2.82).
% Required in the preamble:
%   \usepackage{booktabs}
%   \usepackage{tabularx}
%   \usepackage{svg}        % or convert the SVG to PDF and use \includegraphics
%   \newcolumntype{L}{>{\raggedright\arraybackslash}X}
% Citations refer to fitbuddy_refs.bib.

\section{Application Description}
\label{sec:application}

FitBuddy is a free, open-source (GPL-3.0) health tracker for Android, developed by a single
developer and distributed through GitHub and F-Droid~\cite{fitbuddy-repo,fitbuddy-fdroid}. We
assess version 3.2.82 (commit \texttt{28d1adb}); all statements about the implementation were
verified in its source code. Where the code does not reveal how a component behaves, we state
our assumption explicitly.

\subsection{Functionality}
Users log meals and workouts by taking a photo or typing a free-text description. A large
language model (LLM) estimates the calories and macronutrients, after which the user reviews the
estimate, adjusts ingredient weights and saves it. Packaged foods can be scanned by barcode.
Further features are a daily dashboard, reusable meal presets, structured workouts with an
exercise catalog, body-measurement tracking, progress charts, and an AI ``coach'' that comments on
the user's progress. \emph{Region packs} adapt prompts and staple foods to India, North America,
Europe or Latin America.

During onboarding, the user must connect an AI provider and enter their first and last name, age,
height, weight, optional sex, activity level and goal. Calorie targets are computed on the
device; the AI only picks one of several pre-computed plans.

\subsection{Stakeholders}
\begin{itemize}
  \item \textbf{Users} are the data subjects. They also act as \emph{operators}: they choose the
        AI provider, create the API key and thereby hold the contract with that provider. Users
        who donate become \emph{supporters} and may appear on a public supporter list.
  \item \textbf{The developer} writes and releases the app, but also \emph{operates} the
        server-side services (a Sentry project, a proxy on Vercel and a MongoDB Atlas database)
        and develops the companion app FreeScale. For telemetry and cloud backups, the developer
        determines the purposes and means of processing and therefore acts as
        controller~\cite{gdpr}.
  \item \textbf{AI providers} (OpenRouter, Google Gemini, Ollama, OpenAI or any compatible host)
        receive photos, descriptions and profile data. With a self-hosted Ollama server, the user
        is their own provider.
  \item \textbf{Infrastructure and platform providers}: Sentry, Vercel and MongoDB process data
        on behalf of the developer; Google stores the Android backup; GitHub and F-Droid
        distribute the app and serve update information.
  \item \textbf{Public data services}: Open Food Facts~\cite{off-api} (product data) and the
        jsDelivr CDN (exercise catalog and images).
\end{itemize}
Roles overlap: Google is both a possible AI provider (Gemini) and the backup host, GitHub hosts
both the code and the supporter list, and the developer is involved on both sides of the
FreeScale interface.

\subsection{Collected Data and Purposes}
Table~\ref{tab:data} lists the personal data that FitBuddy processes, why, and where it goes.
Body measurements, dietary and activity records are \emph{data concerning health} and the region
can indicate ethnic or cultural background; both are special categories of personal data under
Art.~9 GDPR~\cite{gdpr}.

\begin{table}[htbp]
\centering
\footnotesize
\caption{Personal data processed by FitBuddy. Flow numbers refer to Figure~\ref{fig:architecture}.}
\label{tab:data}
\begin{tabularx}{\textwidth}{@{}L L l@{}}
\toprule
\textbf{Data} & \textbf{Purpose} & \textbf{Leaves the device via} \\
\midrule
First and last name & Personalised greetings; usage statistics for the developer & 2, 5 \\
Age, sex, height, weight, activity level, goal & Calorie targets; context for AI estimates & 1, 2, 6, 7 \\
Body measurements (weight, body fat, BMR, muscle mass, further smart-scale values) & Progress
  tracking & 1, 2, 6--8 \\
Meals, ingredients, scanned barcodes, presets & Core logging function & 2, 3, 6, 7, 9 \\
Meal photos and free-text descriptions & AI estimation; not stored on the device & 1 \\
Workouts and exercise logs & Core logging function & 2, 6, 7, 9 \\
Region & Adapting prompts and foods; revealed by the regional prompts & 1, 2 \\
Support ID (random UUID) & Identifies cloud backups and crash reports & 5, 6 \\
Device model, name and identifier & Crash diagnosis and usage statistics; labelling backups & 5, 6 \\
AI API keys & Authenticating to the chosen AI provider & 1, 2, 6, 7 \\
\bottomrule
\end{tabularx}
\end{table}

FitBuddy requests no location, contacts or storage permissions; the only sensitive permission is
the camera. Photos are re-encoded before upload, which removes EXIF metadata such as GPS
coordinates.

\subsection{Architecture}
Figure~\ref{fig:architecture} shows the architecture at an abstract level. The app is written in
Kotlin and follows a layered design: a user interface, application logic, and local storage
consisting of a SQLite database (via Room) for health records and a key--value store (Android
DataStore) for settings, identifiers and API keys. There is no user account and no central
application server. All communication with external parties is summarised as flows (1)--(9).

\begin{figure}[htbp]
  \centering
  \includesvg[width=\textwidth]{fitbuddy_architecture}
  \caption{Abstract architecture of FitBuddy and its data flows: (1) AI estimation, (2) progress
  insights and coach chat, (3) barcode lookup, (4) catalog and update requests, (5) telemetry,
  (6) cloud backup, (7) Android Auto Backup, (8) measurement exchange with FreeScale, (9) manual
  export.}
  \label{fig:architecture}
\end{figure}

\subsection{Design and Implementation Choices}
\paragraph{Local-first storage.} All records stay on the device and can be deleted individually;
there is no option to delete all data at once. The app adds no encryption of its own and relies
on Android's file-based encryption.

\paragraph{Bring your own AI provider.} Each AI request (flow 1) contains the photo or text plus a
summary of the profile, today's calorie balance and the latest body measurement. Progress
insights (flow 2) add 30 days of detailed history, monthly averages of older data and the first
name. The user selects the provider; with a self-hosted Ollama server, the data stays in the
user's own network, which is why unencrypted HTTP is allowed for user-configured servers.
Failover between API keys and models never switches to another provider.

\paragraph{Telemetry.} The app uses Sentry~\cite{sentry-android} for crash reports, configured
without IP addresses and with a filter for secrets. Crash reporting is on by default in the
GitHub build without asking, and off by default in the F-Droid build. Independently of this
setting, a daily and post-update \emph{heartbeat} reports the user's full name, the date of their
first log, the number of records, the AI provider and the device model. When Sentry is
unreachable, for example because of a DNS-based blocker, telemetry is relayed through the
developer's Vercel proxy.

\paragraph{Backups.} Three backup mechanisms exist. Android Auto Backup~\cite{android-autobackup}
copies all app data, including API keys, to the user's Google account (flow 7). An optional
cloud backup (flow 6) uploads an AES-256-GCM-encrypted copy of all data to MongoDB Atlas through
the developer's proxy, automatically at most every 12 hours. Its key is derived with PBKDF2 from
a user-chosen password or, by default, from the support ID; the backup is stored and retrieved
under that same ID and cannot be deleted from the app. Finally, users can export a backup file,
which is unencrypted unless a password is set (flow 9).

\paragraph{Integration with other apps.} The measurement bridge is an exported Android
\texttt{ContentProvider} that lets FreeScale read all measurements and add new ones (flow 8). It
accepts callers based on their package name.

\paragraph{Other.} Barcodes are decoded on the device; only the barcode is sent to Open Food Facts
(flow 3). The OpenRouter sign-in uses OAuth with PKCE~\cite{openrouter-oauth}. The GitHub and
F-Droid builds share the same code and differ in crash-reporting default, update source and
signing key.

\subsection{Assumptions}
\begin{itemize}
  \item The server-side code of the developer's Vercel proxy is not public. We assume it stores
        backups unchanged in MongoDB Atlas and forwards telemetry to Sentry without further
        processing.
  \item We treat the AI providers, Sentry, Google and Open Food Facts as black boxes that process
        the received data according to their own policies; their retention periods are not
        considered.
  \item We consider a typical user in the European Union who installs the GitHub build and uses a
        cloud-hosted AI provider. The companion app FreeScale is out of scope apart from its
        interface with FitBuddy.
\end{itemize}
